\section{Introducción}
\label{sec:introduccion}

Este trabajo establece una generación coinductiva común de
\(\pi,\varphi,e\), la clausura dodecafásica de \(\alpha\) y su
representación analítica correlacionada, y la composición algebraica de
la coordenada electrónica. El bloque electrónico determina además una
representación de espín \(1/2\), cuyos proyectores dan las helicidades
\(\pm1/2\) al fijar una dirección de propagación. La misma genealogía
conserva un registro de incidencia que conduce a los diseños de Witt,
los códigos de Golay, el retículo de Leech y \(V^\natural\). El
resultado une, por tanto, valores determinados, memoria recuperable y
realizaciones algebraicas mediante operaciones explícitas.

La determinación de los valores de las constantes fundamentales constituye una
cuestión central de la física matemática. Una teoría puede establecer relaciones
entre observables, identificar simetrías y organizar predicciones de gran precisión,
mientras conserva parámetros cuyo valor se determina experimentalmente.
Comprender por qué esos parámetros adoptan sus valores requiere una construcción
que identifique las operaciones, las normalizaciones y las condiciones de
compatibilidad de las que resultan. Ésta es la cuestión que organiza el presente
trabajo: la relación entre la estructura que determina una constante, la ecuación
que expresa esa determinación y el valor obtenido al evaluarla.

Se aborda esta cuestión mediante un núcleo formal matemático denominado
\emph{Holografía Modular Triádica}, en adelante HMT. Sus tres componentes son
\emph{All Possible Paths}, en adelante APP; TRIT, denominación de la estructura
ternaria orientada; y \emph{Topological Phase Kernel}, en adelante TPK.
APP es el soporte toroidal discreto \(X_9=(\mathbb Z/9\mathbb Z)^2\).
Sus traslaciones cardinales definen el grafo de Cayley de la estructura
aditiva, con generadores \(\pm(1,0),\pm(0,1)\); sobre sus vértices
se calculan las evaluaciones suma y producto, conservando residuos y
cocientes. TRIT es la firma orientada \(\{-1,0,+1\}\), que distingue
los regímenes hiperbólico, parabólico y elíptico. TPK compone la selección
de fase, el transporte de las evaluaciones y la actualización de su
memoria. El grafo, las evaluaciones y la dinámica tienen así funciones
diferentes sobre un mismo estado.

La construcción determina la realización numérica junto con la estructura
que la produce. Las regiones de clausura, propagación y
autoescala se construyen en un dominio finito de condiciones iniciales
orientadas; su prolongación compatible determina posteriormente las coordenadas
\(\pi,e,\varphi\). El registro dodecafásico interviene en la determinación de
\(\alpha\). La construcción electrónica compone esas coordenadas con una
estructura central y con sus operadores de retorno. Finalmente, la incidencia
retenida durante el recorrido admite una realización mediante códigos,
retículos y un álgebra de vértices. Los valores y las relaciones de
incidencia son realizaciones de datos comunes; la conservación de esos
datos durante el refinamiento establece su conexión matemática.

La cuestión afecta al estatuto explicativo de las constantes. Cuando una
teoría recibe un parámetro determinado experimentalmente, sus ecuaciones
establecen las consecuencias que se siguen de ese valor. Una derivación
estructural modifica esa relación: el valor pasa a formar parte de las
consecuencias de la construcción. La pregunta por la intensidad de un
acoplamiento o por una razón característica de la materia se convierte
entonces en la pregunta por las relaciones matemáticas que las determinan.
El tránsito del parámetro a la consecuencia constituye el argumento rector
de este artículo.

La determinación numérica adquiere, desde esta perspectiva, un contenido
constitutivo. La ecuación reúne operaciones cuya procedencia debe hacerse
explícita; su evaluación obtiene el valor que esas operaciones fijan. Una
vez establecidas las condiciones de existencia, unicidad y compatibilidad
correspondientes, la necesidad del resultado pertenece a la demostración.
Explicar por qué una constante adopta un valor significa reconstruir esa
necesidad dentro del dominio de la teoría. La precisión decimal muestra
una consecuencia cuantitativa de la estructura, mientras la derivación
expone la razón por la que se obtiene precisamente esa consecuencia.

El alcance de esta cuestión aumenta cuando varias determinaciones
comparten un mismo origen. La tesis propuesta por los autores vincula
las constantes con la organización que produce sus relaciones recíprocas.
La arquitectura común debe explicar tanto la individualidad de cada
coordenada como las dependencias que la enlazan con las restantes.
El orden de las construcciones adquiere por ello valor demostrativo:
permite reconocer qué datos se generan, qué operaciones los utilizan y
qué relaciones sobreviven a las sucesivas realizaciones. La diversidad
de valores se estudia como expresión de una unidad formal articulada.

En esta organización, la metrología desempeña una función posterior y
esencial. La construcción fija sus resultados antes de confrontarlos con
las magnitudes observadas; la experiencia examina su correspondencia
física y delimita el alcance del acuerdo empírico. La necesidad matemática se establece
respecto de las premisas y operaciones declaradas, y la realización
natural se examina mediante la identificación de observables y su
contrastación independiente. Ambas tareas contribuyen a una explicación
completa: la primera determina el resultado dentro del corpus; la segunda
estudia su pertinencia para describir la constitución cuantitativa del
mundo.

Esta perspectiva permite situar históricamente la investigación. Los
antecedentes considerados a continuación precisan distintas formas de
relacionar leyes, parámetros y observación. El presente trabajo adopta
como problema central la determinación interna de los valores y de la
estructura que los vincula. Sus desarrollos aritméticos, geométricos y
algebraicos reciben su unidad argumental de esa pregunta. La aspiración
ontológica reside en explicar cómo propiedades físicas aparentemente
independientes pueden corresponder a relaciones necesarias de una misma
organización matemática.


Este problema adquiere especial nitidez para los acoplamientos adimensionales.
Una coordenada expresada en unidades depende de la convención metrológica;
una relación adimensional permite preguntar por su determinación con
independencia de esa elección. La constante de estructura fina interviene
en relaciones cuantitativas contrastables y suscita, a la vez, la pregunta
por la estructura que fija el acoplamiento electromagnético. El programa
histórico de su explicación proporciona un antecedente preciso para
estudiar conjuntamente el valor de una constante y la organización
matemática de la que procede.

Un episodio esclarecedor aparece en el trabajo de Dirac de 1931 sobre
singularidades cuantizadas del campo electromagnético. Su propósito se
dirige expresamente a la existencia de una carga eléctrica mínima. El
resultado establece una relación entre esa carga y la intensidad de un
monopolo magnético; Dirac advierte, sin embargo, que la construcción no
determina por sí sola el cociente adimensional vinculado al valor
aproximado \(137\)\cite{dirac1931}. Se obtiene una restricción estructural
genuina cuya fuerza demostrativa no equivale a la fijación independiente
del acoplamiento. Este precedente muestra que la pregunta por el origen
de las constantes formaba parte de la investigación fundamental y exige
identificar exactamente qué determina cada teoría.

La historia de la electrodinámica cuántica precisa la diferencia entre
esa demanda y el cálculo de efectos observables. En su conferencia Nobel
de 1965, Feynman reconstruyó procedimientos que expresaban desplazamientos
de niveles y procesos de dispersión mediante la masa experimental del
electrón\cite{feynman1965}. La renormalización hizo posible obtener
resultados finitos y contrastables conservando esa referencia. La eficacia
de estas predicciones constituye un logro teórico efectivo; su dependencia
de un parámetro experimental identifica, a la vez, el lugar lógico en el
que una determinación adicional podría ampliar la explicación. La
cuestión constitutiva se añade a la capacidad predictiva y estudia las
condiciones que fijan sus parámetros.

El informe conjunto LEP--SLD de 2006 ofrece una manifestación experimental
de esta arquitectura\cite{lep2006}. Sus autores distinguen los parámetros
libres del modelo de las relaciones que éste impone entre observables.
Mediante correcciones radiativas, los datos de la resonancia Z permitían
inferir indirectamente las masas del quark top y del bosón W y contrastarlas
con mediciones directas. Las predicciones entrelazadas y la determinación
experimental de parámetros formaban un mismo procedimiento de prueba.
El problema relevante es, por tanto, identificar las dependencias que una
construcción determina internamente y aquellas cuya especificación sigue
requiriendo información empírica.

En este contexto, se entenderá por \emph{determinación epistemológica}
el procedimiento que establece un valor mediante la relación entre teoría,
observación e inferencia; por \emph{determinación constitutiva}, la
deducción de ese valor a partir de las relaciones que definen el objeto
considerado. Esta segunda expresión designa aquí una exigencia matemática:
explicitar el dominio, los operadores, las normalizaciones y las
compatibilidades que seleccionan el resultado. Su relevancia física
procede de la confrontación posterior con las magnitudes observadas.
El cambio de perspectiva afecta al orden de la explicación: el valor se
examina como consecuencia de una estructura determinada. La aspiración
ontológica consiste en relacionar esa estructura con la constitución del
objeto físico; la prueba matemática y su contraste experimental establecen
el alcance con el que puede sostenerse tal relación.



La constante de estructura fina hace especialmente visible el alcance del
problema. En la descripción electromagnética usual, \(\alpha\) es un
acoplamiento adimensional: expresa una relación física que permanece al
cambiar las unidades. Pauli vinculó la explicación de su valor con la
comprensión de la estructura atomística de la electricidad\cite{pauli1946}.
Aquí se presenta primero la construcción matemática del acoplamiento y,
después, la del electrón al que se atribuye esa interacción. Su proximidad
en el argumento responde a una dependencia estructural: los datos que
determinan el acoplamiento intervienen también en la composición de la
coordenada electrónica. La comparación física examina al final esas
expresiones previamente fijadas.

La perspectiva histórica comprende asimismo la relación entre acción, fase
y orientación. En su formulación por caminos, Feynman asigna a las
trayectorias una fase dependiente de la acción y obtiene la amplitud mediante
la composición de sus contribuciones\cite{feynman1948}. La denominación
\emph{All Possible Paths} recoge esta referencia conceptual. En la construcción
presente, los caminos se definen sobre una estructura aritmética finita;
las reglas de admisibilidad y sus evaluaciones se especifican mediante
operaciones enteras. La relación con una acción física se formula en la
etapa dimensional, una vez construido ese dominio.

El vínculo entre acción y fase angular también aparece en el tratamiento
lagrangiano de Dirac\cite{dirac1933}. La expresión \(\hbar=h/(2\pi)\)
sitúa la acción respecto de una coordenada angular en radianes. El desarrollo
actual estudia previamente los determinantes, cocientes y transportes que
producen una escala adimensional de acción. Su aplicación a una unidad
física se distingue de las razones en las que esa escala común se cancela.
Esta separación resulta decisiva para explicar la composición electrónica
y el papel del retorno cronológico: el cuanto de reloj proporciona una
escala y la sección electrónica conserva su identidad estructural.

En su conferencia Nobel de 1965, Feynman recordó una conversación con
John Archibald Wheeler acerca de la identidad de carga y masa de los
electrones\cite{feynman1965}. Wheeler imaginaba una única línea de universo
que, al recorrer el tiempo con orientaciones alternas, aparecería en una
sección temporal como múltiples electrones y positrones. Feynman distinguió
esta conjetura del electrón único de la representación que incorporó a
su trabajo: un positrón como tramo de línea electrónica con orientación
temporal inversa. Su artículo de 1949 desarrolla esta representación y
reconoce el antecedente de Stückelberg\cite{feynman1949}. Este episodio
sitúa históricamente el problema de la orientación de trayectorias. En HMT,
la bidireccionalidad se expresa mediante involuciones, transportes y
condiciones de retorno cuyo dominio se especifica en cada construcción.

El contenido aritmético comienza con la coexistencia de dos evaluaciones.
Cada fila de la tabla multiplicativa se obtiene por acumulación aditiva;
la reducción modular registra sus residuos, mientras los cocientes conservan
los múltiplos de la base que atraviesan la frontera. La diferencia entre
suma y producto adquiere así una expresión exacta sobre un mismo soporte.
Las composiciones mixtas pueden ser sensibles al orden y originar una
curvatura discreta. TRIT distingue los regímenes elíptico, parabólico e
hiperbólico mediante sus relaciones cuadráticas, y la exponencial proporciona
una ley común a los tres. Esta organización permite estudiar un número
como realización escalar de una sección compatible que conserva más datos
que su coordenada real aislada.

El refinamiento hace explícita esta distinción. Una sucesión compatible de
cilindros de diámetro decreciente determina una coordenada arquimediana;
la trayectoria conserva simultáneamente orientación, cocientes de transporte,
régimen y memoria. La holonomía nonádica identifica el retorno de fase,
mientras su elevación registra el avance de profundidad. La representación
esturmiana describe el orden aperiódico de esa prolongación y proporciona
una formulación operatoria de sus dos fases. El alcance de esta construcción
simbólica se establece por sus operadores y sus invariantes.

El registro \(K\) recibe una definición propia antes de intervenir en
\(\alpha\). Se construye por una transformación integral reversible del
registro firmado; su realización racional \(\kappa\) conserva las doce
componentes, el origen y la orientación fijados. Las diferencias cíclicas
determinan las relaciones sectoriales, y la componente uniforme completa
su reconstrucción. Para \(\alpha\) se distinguen la determinación posicional
por normalización entera y la determinación analítica mediante el resolvente
de vacancias. Su identificación exige conservar la compatibilidad de las
prolongaciones. Las consecuencias sobre \(e+\pi\) se expresarán en términos
de esas relaciones exactas y de las condiciones aritméticas que efectivamente
permiten decidir.

La selección única de \(K\) se establece mediante operaciones sobre
poblaciones explícitas. La sucesión
\(40320\to21902\to19446\to79\to1\) pasa de ordenaciones a
incidencias indexadas con la rejilla, numeradores distintos, candidatos de
la cara excepcional y supervivientes heterotípicos. El censo regional
\(196\cdot197\cdot196\to331\to2\to1\) aplica condiciones de Gram,
ocupación, distancia y orientación a bloques de tres catálogos.
El apartado~\ref{act21:sec:relacion-censos} distingue los dominios de
ambos recuentos; las demostraciones de selección y de inversión precisan
cómo sus salidas se componen con el registro reversible.

La sección electrónica desarrolla el origen de cada factor de la ecuación:
la norma matricial \(\sqrt3/4\), la transferencia regional
\(\varphi/\pi^2\), el término cúbico en \(\alpha\), la normalización
torsional y la memoria multiplicativa de retorno. Dos lectores de la
trayectoria central determinan el mismo triple de coeficientes mediante
operaciones distintas. La misma recurrencia determina un soporte de
\(27\) celdas y dos historias orientadas que la coordenatización integral distingue.
Su lectura armónica identifica tres modos del registro central y expresa
los coeficientes del retorno mediante sus potencias. La ecuación electrónica
conserva así una dependencia explícita de la trayectoria producida.
La estructura local y la coordenada completa se
mantienen durante la realización dimensional. El contraste se presenta
con los ajustes CODATA 2018 y 2022\cite{codata2018,codata2022}, mostrando
los valores centrales, sus incertidumbres y las diferencias calculadas.

La prolongación excepcional retoma los datos de frontera de ese mismo
recorrido. Los cambios de representación pasan de secuencias ternarias a
códigos, de soportes a diseños de incidencia y de clases discriminantes a
retículos. La construcción del vecino de Leech permite aplicar el
procedimiento de Frenkel, Lepowsky y Meurman para obtener \(V^\natural\);
el Monstruo actúa como grupo de automorfismos de esa álgebra, y Moonshine
describe sus trazas graduadas\cite{flm,borcherds}. La realización escalar
y la realización incidencial constituyen las dos proyecciones examinadas
en el artículo: ambas parten de datos identificados y conservan las
relaciones pertinentes de la construcción.

% BEGIN S27_FRAMING
Dos desarrollos precisan la articulación entre los registros enteros
y sus incidencias. El apartado~\ref{carry27:seccion} deriva el cociclo
de acarreo de la división euclídea de las columnas regionales. La tabla
completa de \(27\) entradas, las representaciones polinómicas y las
pruebas de rango establecen la contribución algebraica de los cocientes.
Las identidades de transporte conservan esa contribución junto con el
residuo; las cotas del cilindro determinan la compatibilidad de la
lectura decimal.

El apartado~\ref{star27:comparacion} compara después las estrellas de
\(B_K=H_e\cap P_\pi\), \(B_\varphi=H_\varphi\cap P_\pi\) y
\(B_{\rm APP}=H_{\rm APP}\cap P_\pi\), polarizadas por el soporte
negativo de la cocadena anterior al acarreo. Sus doce completados,
el censo de las \(495\) cuaternas y los estabilizadores del marco
ordenado especifican la información retenida por esa lectura
incidencial. La conexión entre ambos desarrollos reside en los
registros regionales y orientados producidos por APP--TRIT--TPK:
la evaluación entera conserva magnitudes y acarreos, y la realización
de frontera conserva soportes, intersecciones y simetrías.

% END S27_FRAMING


% BEGIN R27_FRAMING_INTRO_ES
La construcción finita permite precisar qué información se transmite en
cada una de estas etapas. El alfabeto de una ventana consta de \(42\)
bloques, obtenidos de las dos firmas del emisor; las seis ventanas producen
el catálogo sobre el que se realizan los empalmes regionales. El ejemplo
de dos registros con el mismo histograma y distinta correlación temporal
identifica una contribución del orden que la lectura marginal omite. Los
cuatro calibres del observable multiplicativo y los controles de clausura
se examinan sobre dominios fijados, conservando la procedencia del registro
y sus condiciones de frontera.

La misma disciplina distingue tres extensiones de APP: variación del
módulo, repetición periódica del soporte y aumento del número de
coordenadas. La compresión multiplicativa en tres clases conserva una
subretícula de índice dos cuya componente de rango dos realiza una
iteración de Pell. La codificación de pares de trits transporta la suma
mediante su acarreo. Posteriormente, la correlación cuadrática construye
la identidad de la coordenatización de seis posiciones; la neutralidad dual reduce
a dos las ocho fronteras saturadas y la orientación fijada selecciona su
representante. El resultado de divisibilidad de nivel tres se sitúa al
término de la realización excepcional, donde su demostración utiliza
la identidad racional ya establecida.
% END R27_FRAMING_INTRO_ES

% BEGIN R28_FRAMING_INTRO_ES
El enlace regional se desarrolla además sobre el multigrafo etiquetado
del emisor y su cociente de fase. Cinco recorridos de longitud mínima
ocho reúnen, en cada región de clausura, los anclajes APP, de propagación
y de autoescala. Sus testigos orientados y su prueba finita de minimalidad
completan la descripción de los empalmes y de la información que conserva
cada representación.
% END R28_FRAMING_INTRO_ES


La enumeración de las condiciones iniciales delimita los estados admisibles;
los certificados de cálculo permiten comprobar sus relaciones finitas y
los registros de procedencia identifican los datos de cada operación.
Las definiciones, los lemas y sus demostraciones establecen la continuidad
entre estos objetos. Los pasos al límite conservan sus hipótesis y sus
mapas de restricción.
La construcción de la sección~\ref{subsec:alpha-publicaciones-completas}
especifica todos los coeficientes de la representación analítica correlacionada
y demuestra su compatibilidad con la clausura entera. La igualdad de los
prefijos de ambas representaciones a toda profundidad se establece en el
teorema~\ref{thm:alpha-dos-publicaciones-completas}, dentro de la clase
geométrica y del dominio allí definidos.
La cuestión rectora permanece en todo
el recorrido: qué estructura determina el valor y qué relaciones subsisten
al cambiar su representación.
